\documentclass[aspectratio=169,11pt]{beamer}
\usepackage[utf8]{inputenc}
\usepackage[T1]{fontenc}
\usepackage[french]{babel}
\usepackage{helvet}
\renewcommand{\familydefault}{\sfdefault}
\usepackage{hyperref}
\definecolor{bleu}{HTML}{0000FF}
\definecolor{fond}{HTML}{F4F4EF}
\setbeamercolor{background canvas}{bg=fond}
\setbeamercolor{structure}{fg=bleu}
\setbeamercolor{title}{fg=bleu}
\setbeamercolor{frametitle}{fg=bleu,bg=fond}
\setbeamercolor{normal text}{fg=black}
\setbeamertemplate{navigation symbols}{}
\setbeamertemplate{itemize item}{\textcolor{bleu}{\tiny$\blacksquare$}}
\setbeamertemplate{footline}{\hfill\footnotesize\insertframenumber\hspace{1em}\vskip6pt}
\hypersetup{colorlinks=true,urlcolor=bleu,linkcolor=bleu}
\title{Normes et examens par les pairs}
\subtitle{Organisation de coopération et de développement économiques}
\author{Veille institutionnelle}
\date{2026-11-03}

\begin{document}
\begin{frame}\titlepage\end{frame}

\begin{frame}{Normes et examens par les pairs (1/2)}
\begin{itemize}
\item Instruments juridiques : décisions, recommandations, déclarations, accords, plus de 250 (à vérifier)
\item Convention contre la corruption d'agents publics étrangers de 1997
\item Principes directeurs pour les entreprises multinationales
\item Convention fiscale modèle, norme commune de déclaration (CRS), plan BEPS et impôt minimum mondial de 15 \%
\end{itemize}
\end{frame}

\begin{frame}{Normes et examens par les pairs (2/2)}
\begin{itemize}
\item Principes sur l'intelligence artificielle de 2019
\item Examens par les pairs : études économiques par pays, examens de l'aide du CAD, politiques sectorielles
\end{itemize}
\end{frame}

\begin{frame}{Normes et examens par les pairs : ressources à suivre}
\small
Veille hebdomadaire.
\begin{itemize}
\item \href{https://legalinstruments.oecd.org}{Instruments juridiques de l'OCDE}\\{\footnotesize Base juridique, publication : mensuelle}
\item \href{https://www.oecd.org/en/topics/policy-areas/tax.html}{Fiscalité et BEPS}\\{\footnotesize Page thématique, publication : hebdomadaire}
\end{itemize}
\end{frame}

\end{document}
